\scanpage{55}

\[
\deleted{\unsure{A}}\kappa\deleted{\unsure{\varphi}}(U)=1_{\mathcal M}\otimes U
\]
Alors \deleted{\unsure{$\kappa$}} $\lambda^*(U)=\kappa(U)$
$\forall U\in$

\margin{\deleted{\unsure{commute aux produits}}\quad\unsure{$\Longrightarrow$}}
\margin{$\Rightarrow\mathcal M$ assez de proj.\quad\deleted{\unsure{$a)$}}}

\margin{\unsure{b)} $1_{\mathcal M}$ est fortement plat
(i.e $U\mapsto1_{\mathcal M}\otimes U$ transforme mono en mono, et comm. aux produits),
un adjoint à gauche $\varphi$
\par\deleted{b) $1_{\mathcal M}$ est projectif \inserted{est projectif} \unsure{de $Ab_k$}.}
\par\unsure{c)} $\varphi$ admet un foncteur résolvant $\lambda_!$ qui transforme projectifs en \unsure{idem}
\par c) $\lambda_!$ admet un adjoint \inserted{$\lambda^*$} (donc ex.)
\par d) $\lambda_!\lambda^*(k)\longrightarrow k$\quad\deleted{\unsure{$Ab_k$}}
}

\[
\operatorname{Hom}_{\mathcal M}(x,\lambda^*(U))\simeq
\operatorname{Hom}(x,1_{\mathcal M}\otimes U)\simeq
\operatorname{Hom}_k(\varphi(x),U)\simeq
\operatorname{Hom}_{\mathrm{Ch}_\bullet(Ab_k)}(\lambda_!(x),U)
\]

\deleted{$\lambda_!\lambda^*(U)$}
\[
\lambda_!(1_{\mathcal M}\otimes_k U)\simeq\lambda_!(1_{\mathcal M})\otimes_k U\longrightarrow U
\]
\margin{provient de
$\lambda_!(1_{\mathcal M})\Rightarrow\lambda_!\lambda^*(1_k)\longrightarrow k$
en tensorisant avec $U$}

Or $\lambda_!\lambda^*k$

\figure{\begin{tikzpicture}
\draw (0,0)--(4,0) .. controls (4.3,0) and (4.6,.1) .. (4.7,.2);
\draw (4.9,.3) .. controls (5,.6) and (4.8,1.3) .. (5,1.4)--(12,1.4);
\draw (8.7,3.2) .. controls (9.7,3.2) and (9.4,2.2) .. (10.5,2.2) .. controls (11.7,2.2) and (11.7,2) .. (12.2,1.6);
\draw (-.4,-.05)--(-.3,-.05);
\end{tikzpicture}}

$\mathcal M$ cat. \inserted{$k$-}abélienne avec $\otimes$ \inserted{$k$-biadd.}
ass. comm. unit., unité $\unsure{1_k|1_{\mathcal M}}$, si

\deleted{\unsure{$1_k$}} a) $\unsure{1_k|1_{\mathcal M}}$ est fortement plat,
i.e. $\kappa:U\longmapsto U\otimes_k\unsure{1_k|1_{\mathcal M}}$ commute aux
$\underset{\longleftarrow}{\lim}$
[i.e transforme mono en mono + commute aux produits]

\deleted{de sorte} i.e $\kappa$ admet un adjoint à gauche
$\exists\varphi:\mathcal M\longrightarrow Ab_k$

(ii) $\varphi$ admet un foncteur résolvant $\lambda_!$
\deleted{qui est} qui admet un adjoint à droite $\lambda^*$
et $\lambda_!$ exact (à g.), et transforme projectifs en projectifs.

(iii) \deleted{\unsure{a)}}
\[
\begin{array}{ll}
\text{\textcircled{1}}&L_i\varphi(1_{\mathcal M})=0\quad\text{pour }i>0\\
\text{\textcircled{2}}&L_0\varphi(1_{\mathcal M})\bigl(=\varphi(1_{\mathcal M})=\varphi\kappa(k)\bigr)\longrightarrow k\quad\text{iso},\\
\text{\textcircled{3}}&\text{les }\lambda_{!i}(1_{\mathcal M})\text{ sont \uline{plats}}
\end{array}
\]
\margin{[\textcircled{1} et \textcircled{3} \deleted{\unsure{$1_{\mathcal M}$}} vrais si $1_{\mathcal M}$ projectif]}

\deleted{\unsure{(iii)}\ $\unsure{\lambda_!\lambda^*k}$}
(iv) $\exists\quad1_{\mathcal M}\xrightarrow\alpha J$ tel que
\[
\left\{\begin{array}{ll}
\text{\textcircled{1}}&\alpha\otimes\mathrm{id}_x\text{ mono.}\quad\forall x\text{ ds }\mathcal M\\
\text{\textcircled{2}}&\deleted{\unsure{\varphi}}L\varphi_\bullet(J\otimes x)\simeq0\quad\text{pour tt }x\text{ dans }\mathcal M
\end{array}\right.
\]

\deleted{\unsure{$\lambda$}} Soit $U$ de $Ab_k$, je dis que
\[
\lambda^*(U)\simeq1_{\mathcal M}\otimes_kU=\kappa(U)
\]

Il suffit de prouver un \unsure{isom'}
\[
\operatorname{Hom}(x,\deleted{\unsure{\lambda^*(U)}}\deleted{\unsure{1_U}}\kappa(U))
\simeq\operatorname{Hom}(\lambda_!(x),U)
\]
\inserted{fonctoriel}
\deleted{pour $x$ projectif} (puis \unsure{inférer}, comme les deux sont exacts à droite en $x$, \unsure{que c'est vrai} pour \unsure{tout|tous} $x$)

\margin{\unsure{$\eta$} adj. de $\varphi,\kappa$}
\[
\operatorname{Hom}(\varphi(x),\deleted{\unsure{\mathrm{Hom}}}U\unsure{?})
\deleted{\simeq\operatorname{Hom}\lambda_!(x)\text{ i.e.}}
\]
\inserted{Car $\varphi(x)=H_0(\lambda_!(x))$}

\[
\lambda_!\lambda^*(U)=\lambda_!(\kappa(U))=\lambda_!(1_{\mathcal M}\otimes_kU)
\simeq\lambda_!(1_{\mathcal M})\otimes U\xrightarrow{\uline{\text{iso}}}U
\]
\margin{$\lambda_!$ ex. à \unsure{g.}}
\margin{déduit de $\lambda_!(1_{\mathcal M})=\lambda_!\lambda^*(k)\longrightarrow k$}

OK

Soit maintenant $l_\bullet$ projectif \inserted{et} acyclique,
donc homotopique à $0$. A prouver que $\lambda_!\lambda^*(l_\bullet)$ aussi.

\mspage{54}
